% ============================================================
% APPENDIX B — FAILURE TAXONOMY
% ============================================================
% Complete classification of all 174 zero-score cases
% Classification methodology: objective, reproducible criteria
% ============================================================

\section*{Appendix B: Failure Taxonomy}
\label{app:failure-taxonomy}

Table~\ref{tab:failure-taxonomy} summarizes the complete failure taxonomy
from all 174 zero-score cases.

\begin{table*}[h!]
\centering
\caption{Failure Taxonomy for 174 Zero-Score Cases (Complete Classification)}
\label{tab:failure-taxonomy}
\begin{tabular}{lrrp{9cm}}
\toprule
\textbf{Category} & \textbf{Count} & \textbf{Percent} & \textbf{Definition / Classification Criteria} \\
\midrule
CONTENT\_GAP & 84 & 48.3\% &
Response exists but missing rubric-required clinical details.
Default category for failures not matching other criteria. \\
SAFETY\_NEUTRALIZATION & 44 & 25.3\% &
Three or more safety corrections applied (diagnostics count threshold),
potentially removing rubric-critical content. \\
MIS\_INTERPRETATION & 32 & 18.4\% &
Grader justification contains indicators of misreading: ``misread,'' ``misunderstood,''
``incorrectly,'' ``off-topic,'' ``does not address,'' or ``fails to address.'' \\
OVERLONG\_UNDER\_SPECIFIC & 12 & 6.9\% &
Very long response ($>$6000 chars) but meets $<$30\% of rubric criteria.
Verbosity obscures key points. \\
STRUCTURE\_MISS & 2 & 1.1\% &
Response too short ($<$500 characters) to contain required rubric items.
Council output compressed or truncated. \\
\bottomrule
\end{tabular}
\end{table*}

\subsection*{B.1 Classification Methodology}

All 174 zero-score cases were classified using \textbf{objective, reproducible criteria}:

\begin{enumerate}
    \item \textbf{Response length}: Measured from \texttt{response\_text} field
    \item \textbf{Rubric breakdown}: Criteria met/missed from \texttt{grade.rubric\_breakdown}
    \item \textbf{Safety corrections}: Count from diagnostics \texttt{safety\_corrections\_applied}
    \item \textbf{Justification patterns}: Keyword matching in \texttt{grade.justification}
\end{enumerate}

Classification rules (applied in order):
\begin{verbatim}
IF safety_corrections >= 3: SAFETY_NEUTRALIZATION
ELIF response_length > 6000 AND criteria_met < 30%: OVERLONG_UNDER_SPECIFIC
ELIF response_length < 500: STRUCTURE_MISS
ELIF justification contains misread indicators: MIS_INTERPRETATION
ELSE: CONTENT_GAP
\end{verbatim}

\subsection*{B.2 Key Findings}

\begin{itemize}
    \item \textbf{CONTENT\_GAP dominates (48.3\%)}: Most failures are responses that exist
    and are reasonably long but miss specific clinical details the rubric requires.

    \item \textbf{SAFETY\_NEUTRALIZATION is the second most common category (25.3\%)}:
    Cases where three or more safety corrections were applied, potentially removing
    rubric-critical content. This is a significant contributor to zero scores.

    \item \textbf{MIS\_INTERPRETATION is significant (18.4\%)}: A substantial fraction
    of failures involve the council misreading or misunderstanding the clinical scenario.
    
    \item \textbf{Length $\neq$ Quality}: Both very short (STRUCTURE\_MISS) and very long
    (OVERLONG\_UNDER\_SPECIFIC) responses fail, indicating rubric alignment is about
    specificity, not verbosity.
\end{itemize}

\subsection*{B.3 Verification}

Full classification data available at: \texttt{docs/research/GPT52\_FAILURE\_TAXONOMY\_COMPLETE.json}

